% Chapter 2 – included by thesis.tex / Submission.tex
\chapter{State of the Art}
\label{ch:sota}
% ==============================================================================

This chapter gives an overview of gaze estimation with ordinary RGB cameras.
First, I explain the basics (\cref{sec:grundlagen}) and the special problems of
webcams (\cref{sec:webcam}). Then I present important datasets
(\cref{sec:datensaetze}), network architectures (\cref{sec:architekturen}) and
pretrained models (\cref{sec:pretrained}). After that, the chapter covers
personalization, domain shift, test-time adaptation and data quality. At the
end, I derive the research gaps that this thesis addresses.

% ------------------------------------------------------------------------------
\section{Basics of Appearance-Based Gaze Estimation}
\label{sec:grundlagen}
% ------------------------------------------------------------------------------
Eye tracking uses images of the eyes or the face to find out where a person is
looking. There are two kinds of output: the three-dimensional gaze direction
(usually two angles, pitch and yaw) or the two-dimensional \emph{point of
regard} (PoR). The PoR is the point on the screen that the person looks at.
This thesis estimates the 2D PoR from images of an ordinary webcam.

\subsection{Classical Methods}
\label{subsec:klassisch}
Appearance-based methods learn a direct mapping from the image to the gaze
target. They do not need a geometric model of the eye. Early methods used
hand-crafted image features, for example Local Binary Patterns (LBP),
Histograms of Oriented Gradients (HoG) or SIFT. These features were combined
with simple regressors such as linear regression, random forests or support
vector regression (SVR). Such methods are fast. With head movement and
changing light, however, they become much less accurate, because the features
do not separate these changes well from the gaze information.

\subsection{First Deep Learning Methods}
\label{subsec:dl_pioniere}
Convolutional neural networks (CNNs) can map face and eye images directly to
screen positions. Krafka et al.~\cite{krafka2016itracker} introduced
\emph{iTracker}, one of the first CNNs for mobile eye tracking without extra
hardware. The model processes four inputs in separate branches: the left eye,
the right eye, the face and a binary \emph{face grid} (where the face is in the
image). Then the features are combined. Without calibration, the error is
\SI{1.71}{\centi\metre} on smartphones and \SI{2.53}{\centi\metre} on tablets.
For personalization, the authors train one SVR per user on the features of
layer \code{fc1}. They use recordings of up to 13 fixed calibration points.
With all 13 points, the error drops to \SI{1.34}{\centi\metre} and
\SI{2.12}{\centi\metre}. With very few points, however, the SVR easily runs into
\emph{overfitting}: the model almost memorizes the few training samples and
therefore becomes worse on new images. The model runs in real time on mobile devices.

An important contribution of this work is the GazeCapture dataset
(\cref{sec:datensaetze}). It was collected by \emph{crowdsourcing} on smartphones and
tablets in everyday conditions (\enquote{in the wild}): many volunteers
recorded the data themselves with an app on their own device. It is still used as a
benchmark and as a source for pretraining.

\paragraph{Relevance for this thesis.}
iTracker introduced three ideas that are still important: (1) the GazeCapture
dataset, (2) personalization with calibration data and (3) combining features
from several image regions. This thesis uses the same data. Instead of AlexNet
and SVR, however, it uses a ResNet-18 \cite{he2016resnet} with a
regression head and L1 loss. Here the regression head is a multi-layer
perceptron (MLP), a small network of fully connected layers that computes the screen
coordinates $(x, y)$ from the image features of the ResNet-18. It also evaluates the base models with a
subject-disjoint split (\code{norm\_subject}) and not only with a random
split (\code{norm\_random}). Wrong labels are handled already during recording: the
recording program stores the phase of the target for every frame
(\code{WAIT} or \code{MOVE}).

\subsection{Types of Methods}
\label{subsec:taxonomie}
Gaze estimation methods can be divided into three classes:
\begin{enumerate}
  \item \textbf{Model-based (geometric) methods} use a 3D model of the
        eyeball, pupil ellipses and light reflections on the cornea. They
        usually need high-resolution images and infrared light. In normal
        daylight they often fail.
  \item \textbf{Feature-based methods} find geometric points in the face, for
        example eye corners or the centre of the iris. From these points they
        compute the gaze.
  \item \textbf{Appearance-based methods} use the eye or face image itself as
        input. A deep neural network directly predicts 2D coordinates or 3D
        angles.
\end{enumerate}
Thanks to progress in deep learning, appearance-based methods are now the
standard under \emph{unconstrained conditions}, that is, with free head pose,
changing light and any distance to the camera.

Cheng et al.~\cite{cheng2024review} also divide \emph{remote-camera} methods
(the camera stands in front of the person and is not worn on the head) into
three groups: methods with a reconstructed 3D eye model (the eyeball is
modelled geometrically, and the gaze is computed from the pupil and light
reflections), regression on 2D eye features, and appearance-based methods. They sort deep learning approaches by
four aspects: feature extraction, network architecture, personal calibration,
and devices and platforms. An important contribution is a benchmark. It unifies
the pre-processing (face and eye detection, data normalization) and the
post-processing (conversion between 2D gaze point and 3D gaze vector). This
allows a fair comparison of methods. According to the benchmark,
appearance-based methods reach about $5^\circ$ at short distance or with high
resolution. At large distance and low resolution (for example Gaze360 and
RT-GENE), the errors are much larger. As open topics, the authors name: (1)
more robust features that do not depend on the environment, (2) high accuracy
with fast and easy calibration -- here there is a trade-off between
calibration time and accuracy -- and (3) the explainability of the learned
features.

Other surveys give additional background. Liu et al.~\cite{liu2022beholder}
describe the path from classical methods to deep learning; I use this work
mainly as a historical reference. Pathirana et al.~\cite{pathirana2022survey}
cover deep learning methods for single and multiple users. They compare CNN
architectures, some of which are pretrained on ImageNet. Newer pretrained
models such as DINOv2, CLIP or MAE are not the focus of these surveys. For
benchmark comparisons I therefore mainly use Cheng et
al.~\cite{cheng2024review}.

% ------------------------------------------------------------------------------
\section{Webcam-Based Eye Tracking: Everyday Use versus Laboratory}
\label{sec:webcam}
% ------------------------------------------------------------------------------
Webcams are cheap and almost everywhere: in laptops, on desktop computers and
in smartphones. No extra hardware is needed. But the conditions are much harder
than for special eye trackers with infrared light:

\begin{itemize}
  \item \textbf{Infrared systems} use high-resolution cameras, infrared light
        and reflections on the cornea. Liu et al.~\cite{liu2022beholder} report
        typical accuracies of $0.5^\circ$ to $2^\circ$ for such methods.
        Commercial devices work at fixed distances (for example
        \SIrange{50}{80}{\centi\metre}) and reach $0.3^\circ$ to $0.5^\circ$
        under optimal conditions.
  \item \textbf{Webcams in everyday use} deliver low-resolution images without
        controlled light. The ambient light changes, the position of camera
        and screen is different on every device, and distance and head pose
        are free. Liu et al.~\cite{liu2022beholder} report typical errors above
        $3^\circ$ for appearance-based methods. The benchmark of Cheng et
        al.~\cite{cheng2024review} shows about $5^\circ$ at short distance and
        much more at large distance and low resolution.
\end{itemize}

Two well-known works still show that webcam eye tracking is useful in practice.
Papoutsaki et al.~\cite{papoutsaki2016webgazer} present \emph{WebGazer}, an
open-source eye tracker that runs completely in the browser. WebGazer does not
need its own calibration. With \emph{ridge regression}, it continuously learns from
mouse clicks and mouse movements which eye features belong to which screen
position. Ridge regression is a linear regression with a penalty on large
weights (L2 regularization); this keeps the model stable even with few and
noisy data. The accuracy is enough to find the rough area a person looks at, but
not for detailed analysis. Valliappan et al.~\cite{valliappan2020smartphone},
in contrast, pretrain a CNN on GazeCapture, fine-tune it with calibration data
and add a small regression model per person. On smartphones at
\SIrange{25}{40}{\centi\metre} distance they reach $0.6^\circ$ to $1^\circ$.
About 30~seconds of calibration reduce the error from \SI{1.92}{\centi\metre}
to \SI{0.46}{\centi\metre}. According to the authors, this is as accurate as
much more expensive mobile eye trackers. The conditions, however, were
favourable (head almost frontal, short distance). Both works show: without
special hardware, personalization is the most important factor for accuracy.

Properties of the person also affect the accuracy, for example eye shape, face
shape, glasses and the usual head pose. Generic models therefore often have a
systematic offset for new persons. This is why personalization is a central
research topic (\cref{sec:personalisierung}). This thesis studies exactly such
conditions, among others two devices with different webcams,
glasses and backgrounds.

% ------------------------------------------------------------------------------
\section{Large Datasets for Gaze Estimation}
\label{sec:datensaetze}
% ------------------------------------------------------------------------------
Modern appearance-based methods need large and varied datasets.
\Cref{tab:datensaetze} summarizes the datasets that are important for this
thesis.

\begin{itemize}
  \item \textbf{MPIIGaze}~\cite{zhang2015mpiigaze} is one of the first
        standardized everyday datasets. It contains \num{213659} face images of
        15 persons. They were recorded over several months during normal work
        on a laptop~\cite{zhang2019mpiigaze}. The variant MPIIFaceGaze with
        face images is usually evaluated with leave-one-subject-out (each
        person is the test person once).
  \item \textbf{GazeCapture}~\cite{krafka2016itracker} contains \num{2445504}
        images of \num{1474} persons, recorded by crowdsourcing on iPhones and
        iPads. Light and head pose vary a lot. GazeCapture is the data basis of
        the base models in this thesis. Reasons for this choice: it is one of
        the largest and most diverse datasets (many persons, devices and
        environments), and its labels are 2D points on the screen (point of
        regard) -- exactly the quantity this thesis estimates.
  \item \textbf{ETH-XGaze}~\cite{zhang2020ethxgaze} is a high-resolution
        dataset with several cameras. It contains more than one million images
        of 110 persons with very large head and gaze angles.
  \item \textbf{EVE}~\cite{park2020eve} contains synchronized screen content,
        camera images and gaze data. 54 persons were recorded with four
        cameras, more than 12~million images in total.
  \item \textbf{Gaze360}~\cite{kellnhofer2019gaze360} contains indoor and
        outdoor recordings of 238 persons. The 3D gaze labels cover a very wide
        range of head poses and distances (about \num{172000} images in
        training, validation and test).
  \item \textbf{EyeDiap} and \textbf{RT-GENE} are mainly used for tests across
        datasets. RT-GENE contains, among others, larger camera distances and
        head-mounted devices.
\end{itemize}

\begin{table}[htbp]
  \centering
  \small
  \caption{Overview of selected gaze estimation datasets.}
  \label{tab:datensaetze}
  \begin{tabularx}{\textwidth}{L{2.4cm} L{2.6cm} L{2.2cm} Y}
    \toprule
    \textbf{Dataset} & \textbf{Size} & \textbf{Persons} & \textbf{Recording conditions} \\
    \midrule
    MPIIGaze      & \num{213659} images & 15 & Laptop webcam, everyday use over months \\
    GazeCapture   & \num{2445504} images & \num{1474} & Smartphones and tablets, crowdsourcing \\
    ETH-XGaze     & >1~million images & 110 & Lab with several cameras, extreme head poses \\
    EVE           & >12~million images & 54 & Screen and cameras, video sequences \\
    Gaze360       & approx.\ \num{172000} images & 238 & Indoor and outdoor, many head poses and distances \\
    \bottomrule
  \end{tabularx}
\end{table}

% ------------------------------------------------------------------------------
\section{Network Architectures for Gaze Regression}
\label{sec:architekturen}
% ------------------------------------------------------------------------------
\subsection{Convolutional Neural Networks}
\label{subsec:cnn}
CNNs such as ResNet-18/50~\cite{he2016resnet}, VGG or MobileNet are the usual
feature extractors for eye and face images. Their local convolutions detect
fine textures and edges in the eye region. Special gaze models are also often
based on CNNs: FAZE~\cite{park2019faze}, for example, uses a DenseNet-based
encoder-decoder to separate gaze, head pose and appearance. This thesis uses
ResNet-18.

\subsection{Vision Transformers}
\label{subsec:vit}
Vision Transformers (ViTs)~\cite{dosovitskiy2021vit} split an image into small
patches. With self-attention they capture relations across the whole image,
for example between the face, the head rotation and the eyes.
UniGaze~\cite{qin2026unigaze} shows that ViTs pretrained with MAE (sizes Base,
Large and Huge) are very good backbones for gaze estimation.
TTP-GE~\cite{nagpure2025ttpge}, in contrast, combines two ResNet-18 branches
with a transformer fusion. GazeCLIP~\cite{wang2026gazeclip} shows that
transformers are not always better: there, a CLIP-pretrained ResNet-50 reaches
$3.5^\circ$ on MPIIFaceGaze, while the ViT-B/32 variant reaches only
$3.7^\circ$. Ye et al.~\cite{ye2024transgaze}, on the other hand, show with
\emph{TransGaze} that even plain ViTs with large-scale pretraining achieve very
good results on GazeCapture and MPIIFaceGaze. Self-supervised DINO pretraining
worked best there. One should therefore test CNNs and transformers equally and
not assume that one architecture is generally better.

\subsection{Pretrained Foundation Models}
\label{subsec:foundation}
In transfer learning, one uses a model that was pretrained on a very large
number of images. Its features are useful in general. One can freeze the
backbone and train only a small regression head on top. Four kinds of
pretraining are important for this thesis:
\begin{itemize}
  \item \textbf{Supervised pretraining on ImageNet}: the model learns with
        labels given by humans (here: object classes). This is the classical
        starting point for CNNs such as ResNet~\cite{he2016resnet}. This
        thesis uses such a model.
  \item \textbf{DINOv2}~\cite{oquab2023dinov2}: learns robust features with
        \emph{self-supervised learning}: the learning target comes from the
        images themselves, without manual labels.
  \item \textbf{CLIP}~\cite{radford2021clip}: learns image and text together
        by contrastive learning on image-text pairs.
  \item \textbf{Masked autoencoders (MAE)}~\cite{he2022mae}: learn by
        reconstructing hidden parts of the image.
\end{itemize}
The next section describes how these features are used for gaze estimation.

% ------------------------------------------------------------------------------
\section{Pretrained Models for Gaze Estimation}
\label{sec:pretrained}
% ------------------------------------------------------------------------------
\subsection{Frozen Backbone or Full Fine-Tuning}
\label{subsec:frozen_full}
There are two basic strategies to adapt a pretrained model to gaze estimation.
In \emph{full fine-tuning}, all weights are trained. The model adapts most
strongly, but it needs a lot of data. With little data there is a risk of
overfitting or \emph{catastrophic forgetting}: the model loses what it learned
during pretraining. With a \emph{frozen backbone} (for example ResNet-18 or
ViT-B/14), the pretrained weights stay fixed. Only a small regression head is
trained. This gives far fewer trainable parameters, and the general features
are kept. This is important when only little calibration data is available.

Kumar et al.~\cite{kumar2022finetuning} explain this trade-off in theory and
with experiments. They compare full fine-tuning with \emph{linear probing},
where only the last linear layer is trained. On ten datasets with distribution
shift, fine-tuning is on average 2\,\% better on data like the training data,
but 7\,\% worse on shifted data. The reason: when head and backbone are trained
at the same time and the head is still random at the start, the pretrained
features get distorted. The authors therefore propose LP-FT: first linear
probing, then fine-tuning. This is directly relevant for personalization. The
calibration data of one person are few and cover the later use only partly. A
frozen backbone can protect the generalization here.

The regression head is usually a multi-layer perceptron (MLP). It maps the
features, for example the 512 numbers of a ResNet-18, to the coordinates
$(x, y)$. Important choices are the output activation (for example sigmoid,
Gaussian or ReLU) and the loss function (for example L1, smooth L1 or variants
weighted by position). Between the two extremes there are parameter-efficient
methods such as adapters or LoRA~\cite{hu2022lora} (\cref{subsec:peft}).

\subsection{Large-Scale MAE Pretraining: UniGaze}
\label{subsec:unigaze}
Qin et al.~\cite{qin2026unigaze} present UniGaze, a large pretraining approach
made specially for gaze estimation. With masked autoencoders, the model learns
from about 1.6~million normalized face images of more than \num{260000}
persons. The data come from real sources (VGGFace2, VFHQ, CelebV-Text) and from
synthetic datasets (FaceSynthetics, SFHQ-T2I, FFHQ-NV, XGaze-Dense). A central
result is how important data normalization is. Here, the face image and the
gaze label are rotated and scaled into a common space. This removes a large
part of the variation caused by head pose and distance. Zhang et
al.~\cite{zhang2018normalization} were the first to study this method in
detail. A variant without a scaling factor improves the accuracy by 9.5 to
32.7\,\%, depending on the test. UniGaze confirms this: in the ablations
(UniGaze-B, real pretraining data only, leave-one-dataset-out), normalization
reduces the angular error on XTest from $14.56^\circ$ to $11.95^\circ$. With
additional synthetic data, $11.78^\circ$ is reached.

UniGaze uses ViTs of the sizes Base, Large and Huge. During pretraining, 75\,\%
of the image is hidden. ViT-Huge is pretrained for 300 epochs, which takes
about \SI{120}{\hour} on four NVIDIA H100 GPUs. Then the models are fully
fine-tuned with L1 loss on MPIIFaceGaze, ETH-XGaze, EYEDIAP, Gaze360 and
GazeCapture. The authors test three protocols:
\begin{itemize}
  \item \textbf{Domain generalization} (training on ETH-XGaze or Gaze360, test
        on MPIIFaceGaze or the screen part of EYEDIAP): UniGaze-H reaches
        $5.57^\circ$ (X$\rightarrow$M), $4.65^\circ$ (X$\rightarrow$ECS),
        $5.65^\circ$ (G360$\rightarrow$M) and $5.86^\circ$
        (G360$\rightarrow$ECS). This is better than other methods, for example
        CLIP-Gaze with $6.41^\circ$, $7.51^\circ$, $6.89^\circ$ and
        $7.06^\circ$.
  \item \textbf{Leave-one-dataset-out} (training on four datasets, test on the
        fifth): UniGaze-H reaches $11.29^\circ$ (XTest), $4.51^\circ$
        (MPIIFaceGaze), $5.47^\circ$ (GazeCapture), $5.88^\circ$ (EYEDIAP) and
        $12.37^\circ$ (Gaze360).
  \item \textbf{Joint training on all datasets}: UniGaze-H reaches
        $4.46^\circ$ (XTest), $5.08^\circ$ (MPIIFaceGaze), $3.20^\circ$
        (GazeCapture), $5.16^\circ$ (EYEDIAP) and $9.07^\circ$ (Gaze360).
\end{itemize}

The ablations show: more pretraining data improve generalization, and real plus
synthetic data are better than real data only. Normalization and many
different head poses are also important. If the variety of head poses is
limited, the error on XTest rises from $11.95^\circ$ to $16.70^\circ$.
Comparisons with DINO and FaRL suggest that the composition of the data can be
more important than their amount.

\paragraph{Limitations.}
UniGaze leaves some questions open that matter for personalization with
webcams. The results are based on complete datasets and full fine-tuning, not
on a short calibration with few images of one person. Lightweight strategies
such as head-only training, LoRA or adapters are not the focus. The effect of
wrong labels in short calibrations is not studied either. Pretraining also
needs a lot of computing power (ViT-Huge: about \SI{120}{\hour} on
$4\times$H100). It also remains open whether the improvements hold under real
webcam conditions -- with glasses, changing backgrounds, low resolution and
different webcams.

\paragraph{Relevance for this thesis.}
UniGaze is an important reference for gaze-specific pretraining. UniGaze
studies generalization across many persons. This thesis asks the complementary
question: how well can a pretrained model be adapted to one single person with
a short calibration (here \SIrange{30}{120}{\second}, that is, about 180 to
720 labels)? This thesis uses a ResNet-18 pretrained on ImageNet for this.
UniGaze, DINOv2 and CLIP models are a natural next step
(\cref{tab:unigaze_vergleich}).

\begin{table}[htbp]
  \centering
  \small
  \caption{Comparison of UniGaze and this thesis.}
  \label{tab:unigaze_vergleich}
  \begin{tabularx}{\textwidth}{L{2.8cm} L{3.2cm} L{2.6cm} Y}
    \toprule
    \textbf{Criterion} & \textbf{UniGaze-H} & \textbf{UniGaze-B} & \textbf{This thesis} \\
    \midrule
    Pretraining data & 1.6~million face images, >\num{260000} persons & Same data & ImageNet (ResNet-18), regression head on GazeCapture \\
    Backbone & ViT-Huge & ViT-Base & ResNet-18, frozen \\
    Fine-tuning & Whole model & Whole model & Regression head only \\
    Calibration data & Complete training datasets & Complete training datasets & \SIrange{30}{120}{\second} per session (approx.\ 180--720 labels) \\
    Evaluation & Across datasets, domain generalization & Across datasets, domain generalization & Subject-disjoint base models, personalized webcam sessions \\
    Wrong labels & No phase-based filtering & No phase-based filtering & \code{WAIT}/\code{MOVE} logging \\
    Compute & Approx.\ \SI{120}{\hour} on $4\times$H100 & Only given for ViT-H & No pretraining; personalization approx.\ \SI{2}{\second} on a CPU \\
    Angular error (leave-one-dataset-out) & $4.51^\circ$--$12.37^\circ$ & $4.73^\circ$--$12.41^\circ$ & Diagonal error in \%, not directly comparable \\
    \bottomrule
  \end{tabularx}
\end{table}

\subsection{Image-Text Models: CLIP-Gaze and GazeCLIP}
\label{subsec:clip}
Two recent works use the image-text model CLIP~\cite{radford2021clip} for gaze
estimation, but with different goals.

\paragraph{CLIP-Gaze.}
Yin et al.~\cite{yin2024clipgaze} want to improve domain generalization with
CLIP-Gaze. The model should stay good in a new environment without labelled
data from that environment. The idea: factors that have nothing to do with
gaze are described by texts, such as \enquote{An image of a face with
\{condition\}.} These factors form three groups: appearance (identity, facial
expression, hair, age, etc.), accessories (glasses, hat, make-up, etc.) and
image quality (noise, blur, light). Three losses guide the training: a
distillation loss aligns the gaze features with the CLIP image encoder. A
separation loss makes them different from the gaze-irrelevant features. A rank
loss makes sure that similar gaze directions also have similar features. In
addition, \emph{Personalized Context Optimization} (PCO) adapts the texts to
each person. CLIP stays frozen. A separate network of ResNet-18, an MLP filter
and a regression head is trained. With ETH-XGaze and Gaze360 as source and
MPIIFaceGaze and EyeDiap as target, CLIP-Gaze reaches about $6.97^\circ$ on
average over four tasks.

\paragraph{GazeCLIP.}
Wang et al.~\cite{wang2026gazeclip} use coarse gaze information from text for
supervised gaze regression. Without training (zero-shot), a prompt module uses
texts such as \enquote{A photo of a face gazing [class]} to find a coarse
direction from \{Front, Down, Left, Right\}. The image encoder is a
CLIP-pretrained ResNet-50 and is fine-tuned. The text encoder (a transformer
with 12 layers) stays frozen. A cross-attention module connects both: the text
is the query, and the image features are key and value. An MLP with three
layers ($1024 \rightarrow 256 \rightarrow 128 \rightarrow 2$) predicts pitch and
yaw. GazeCLIP reaches $3.5^\circ$ on MPIIFaceGaze, $4.7^\circ$ on EyeDiap and
$7.3^\circ$ on RT-GENE. Cross-attention ($3.6^\circ$) is better than simple
concatenation ($4.2^\circ$) or addition ($4.4^\circ$). A fine-tuned image
encoder with a frozen text encoder ($3.6^\circ$) is much better than two frozen
encoders ($8.9^\circ$). Errors mainly occur in bad light and when hair or
glasses cover the eyes. In addition, the zero-shot classification can give the
wrong direction in a new environment.

\paragraph{Limitations and relevance for this thesis.}
Both works show that CLIP features are useful for gaze estimation. They do not
study personalization with a short calibration, however. For CLIP-Gaze, it
stays open whether adapting the CLIP image encoder improves the estimate for one
person. GazeCLIP needs full supervised training on benchmark data. It covers
neither personalization with few images nor wrong calibration labels, and
according to the publication its data are only available on request. CLIP
features were not yet studied in this thesis; they are a candidate for the
outlook (\cref{tab:clip_vergleich}). The angular errors of these works are only
reference values. They cannot be compared directly with the diagonal error of
this thesis.

\begin{table}[htbp]
  \centering
  \small
  \caption{Comparison of CLIP-Gaze, GazeCLIP and this thesis.}
  \label{tab:clip_vergleich}
  \begin{tabularx}{\textwidth}{L{2.6cm} Y Y Y}
    \toprule
    \textbf{Aspect} & \textbf{CLIP-Gaze~\cite{yin2024clipgaze}} & \textbf{GazeCLIP~\cite{wang2026gazeclip}} & \textbf{This thesis} \\
    \midrule
    Goal & Domain generalization & Supervised gaze regression & Personalization with webcam and little data \\
    Use of CLIP & Image and text encoder frozen & Image encoder fine-tuned, text encoder frozen & Not used (outlook) \\
    Architecture & ResNet-18, MLP filter, regression head & CLIP ResNet-50, cross-attention, MLP & ResNet-18 (ImageNet), regression head \\
    Labels of the target person & None & Complete training data & \SIrange{30}{120}{\second} calibration \\
    Reported error & Approx.\ $6.97^\circ$ (mean of 4 tasks) & $3.5^\circ$ / $4.7^\circ$ / $7.3^\circ$ & Diagonal error in \% \\
    Data quality & Models gaze-irrelevant factors & Coarse direction from text (uncertain) & \code{WAIT} frames are removed \\
    \bottomrule
  \end{tabularx}
\end{table}

% ------------------------------------------------------------------------------
\section{Data-Efficient Personalization}
\label{sec:personalisierung}
% ------------------------------------------------------------------------------
Generic models often have a systematic offset for new persons. It is caused by
differences in the eye, the face shape and the head pose. Park et
al.~\cite{park2019faze} speak of anatomical differences between persons that
limit the accuracy of generic networks. Personalization tries to close this gap
with as little calibration effort as possible. The following sections describe
classical calibration, few-shot and meta-learning, and parameter-efficient
methods.

\subsection{Calibration Methods}
\label{subsec:kalibrierung}
In active calibration, the person looks at known points on the screen one after
the other. The coordinates of these points are the labels for the adaptation.
Typical are a few fixed points or a point that moves across the screen.
iTracker~\cite{krafka2016itracker} uses up to 13 fixed points and trains one
SVR per user on them (\cref{subsec:dl_pioniere}). He et
al.~\cite{he2019ondevice} show that good methods with only 2 to 5 points per
person are better than earlier methods with more than 13 points. On
smartphones, their best model with at most five points is 24 to 26\,\% more
accurate than the state of the art at that time and runs in real time on the
device. They also propose a method that needs only unlabelled face images.
Valliappan et al.~\cite{valliappan2020smartphone} use, in addition to points, a
\emph{zigzag calibration}: a target moves slowly for 60~seconds from top left
to bottom right while the front camera records at \SI{30}{\hertz}. Already
about 100 images reduce the error almost to one quarter.

The methods differ not only in \emph{how much} data they need, but also in
\emph{which part} of the model they adapt:
\begin{itemize}
  \item \textbf{Offset calibration:} Chen and Shi~\cite{chen2020offset} split
        the gaze angle into a part that is the same for all persons and a
        personal offset (bias). They observe that the offset between persons is
        larger than the variation within one person. The calibration then
        needs only one gaze point. On MPIIGaze and EYEDIAP, the method is up to
        10\,\% better than comparable models. Already about nine calibration
        images reduce the error clearly.
  \item \textbf{Few person parameters:} Lindén et al.~\cite{linden2019spaze}
        describe personal differences with the \emph{Spatial Adaptive Gaze
        Estimator} (SPAZE) by a few latent parameters. Calibration is then only
        a small optimization problem and hardly prone to overfitting. With nine
        examples, SPAZE reaches $2.70^\circ$ on MPIIGaze, 14\,\% better than
        the state of the art at that time. Three parameters per eye are enough
        to describe the personal differences well.
  \item \textbf{Differential estimation:} Liu et
        al.~\cite{liu2021differential} train a CNN that predicts the
        \emph{difference} in gaze between two eye images of the same person.
        Together with labelled reference images from the calibration, this
        gives the gaze in the new image. Because both images come from the same
        person, disturbances such as eye shape, eyelid closure or light partly
        cancel out. On EYEDIAP, MPIIGaze and UT-Multiview, the method beats the
        state of the art with only one calibration image.
  \item \textbf{Personalization without calibration:} Zhao et
        al.~\cite{zhao2022ebpn} propose the \emph{Episode-based Personalization
        Network} (EbPN). A separate branch of the network should compensate
        personal deviations without calibration data.
\end{itemize}
Together, these works show: a large part of the personal error can be explained
with a few adapted parameters. This supports the strategy of this thesis to
adapt only a small regression head on a frozen backbone.

Cheng et al.~\cite{cheng2024review} name a basic trade-off: longer
calibrations make the estimate more accurate, but they are less convenient.
High accuracy with a short and easy calibration is therefore an open research
question. Another point is often overlooked: the quality of the calibration
labels themselves. When a target jumps, the gaze follows only after a short
delay. Frames from this transition can have a wrong label
(\cref{sec:label_noise}).

\subsection{Few-Shot Learning and Meta-Learning: FAZE}
\label{subsec:faze}
Park et al.~\cite{park2019faze} present FAZE (\emph{Few-Shot Adaptive Gaze
Estimation}), an influential method for personalization with few calibration
images. It combines separated (disentangled) features for gaze, head and
appearance with Model-Agnostic Meta-Learning (MAML). FAZE does not train a new
model for every person. Instead, it learns a starting point from which a few
examples are enough for a fast adaptation.

FAZE has two parts. A \emph{Disentangling Transforming Encoder-Decoder}
(DT-ED) based on DenseNet represents appearance, gaze direction and head pose
in separate codes. Because these codes rotate together with the input
(rotation-equivariant), changes of gaze and head can be described as rotations
in feature space. Training uses image pairs of the same person with a
reconstruction loss, a gaze loss and a consistency term. The second part is a
gaze estimator that is meta-trained with MAML: over many training persons it
learns a starting point from which a few gradient steps produce a personal
model. This should avoid overfitting, which is a risk for normal fine-tuning
with very little data.

FAZE is trained on GazeCapture and tested on GazeCapture and MPIIGaze. With
$k = 9$ calibration images, the mean error on MPIIGaze is $3.14^\circ$. For all
$1 \le k \le 256$, FAZE is better than the compared methods. On GazeCapture,
the error drops by 19\,\% compared with the state of the art, from
$3.91^\circ$ to $3.18^\circ$.

\paragraph{Limitations and differences.}
FAZE shows that calibration depends not only on the amount of data, but also on
the features learned before and on the type of adaptation. The features of
FAZE, however, come from supervised training with gaze data and a DenseNet.
General pretrained models (DINOv2, MAE, CLIP) are not studied. The adaptation
uses meta-learning; simpler ways such as head-only training, LoRA or adapters
are not the focus. Webcam-specific problems such as changing image quality,
glasses, backgrounds and different cameras are not considered either, and
neither is the removal of calibration frames by recording phase. This thesis
takes up some of these points. Because FAZE reports angular errors and this
thesis uses an error relative to the screen diagonal, the numbers are not
directly comparable.

\begin{table}[htbp]
  \centering
  \small
  \caption{Comparison of FAZE and this thesis.}
  \label{tab:faze_vergleich}
  \begin{tabularx}{\textwidth}{L{3.0cm} Y Y}
    \toprule
    \textbf{Aspect} & \textbf{FAZE~\cite{park2019faze}} & \textbf{This thesis} \\
    \midrule
    Goal & Personalization of the gaze direction with few images & Data-efficient personalization of the screen position with a webcam \\
    Features & Disentangled codes, DenseNet encoder-decoder & ResNet-18, pretrained on ImageNet and frozen \\
    Pretraining & Supervised on GazeCapture & ImageNet; regression head on GazeCapture \\
    Adaptation & MAML and gradient steps per person & Regression head only \\
    Calibration data & $k = 1$ to $256$ images & \SIrange{30}{120}{\second} (approx.\ 180--720 labels) \\
    Label quality & No phase-based filtering & \code{WAIT} frames removed via phase logging \\
    Evaluation & GazeCapture and MPIIGaze & 18 webcam sessions under changing conditions \\
    Reported error & $3.14^\circ$ (MPIIGaze, $k = 9$) & Diagonal error in \%, not directly comparable \\
    \bottomrule
  \end{tabularx}
\end{table}

Yu et al.~\cite{yu2019redirection} extend adaptation with few images: from a
few reference images of a person they create additional training images by gaze
redirection and use them to adapt a generic model.

\subsection{Parameter-Efficient Fine-Tuning (PEFT)}
\label{subsec:peft}
Parameter-efficient methods adapt only a small part of the parameters. The
pretrained backbone stays almost unchanged. This saves memory and computing
time and lowers the risk of overfitting when only little calibration data is
available. Three variants are important for this thesis:
\begin{itemize}
  \item \textbf{Head-only training (FC only):} The backbone stays completely
        frozen; only the regression head is trained. The backbone features can
        be computed once and stored. This makes the method especially fast.
        This thesis uses it.
  \item \textbf{Adapters}~\cite{houlsby2019adapter}: Small extra modules are
        inserted between the layers of a frozen network. Only they are
        trained.
  \item \textbf{Low-Rank Adaptation (LoRA)}~\cite{hu2022lora}: The pretrained
        weight matrices $W$ stay fixed. Only a change $\Delta W = BA$ made of
        two narrow matrices with low rank $r$ is trained. This gives far fewer
        trainable parameters.
\end{itemize}
In gaze estimation, such methods have hardly been studied systematically. One
exception is TPGaze~\cite{liu2024tpgaze}. It transfers prompt tuning from
language processing: an extra prompt with less than 1\,\% of the parameters of
a ResNet-18 is adapted at test time without labels. Meta-learning makes sure
that this adaptation really lowers the gaze error (\cref{sec:tta}). Bao and
Lu~\cite{bao2024feature2gaze} also show that especially the last fully
connected layer, which maps many image features to the gaze, overfits easily.
So the design of the head matters even for head-only training. The other works
discussed here use full fine-tuning (UniGaze, GazeCLIP), meta-learning (FAZE)
or a frozen model that guides a separately trained network (CLIP-Gaze). This
thesis studies head-only training for personalization with short webcam
calibrations. Adapters and LoRA are left for future work.

% ------------------------------------------------------------------------------
\section{Domain Shift, Generalization and Evaluation Protocols}
\label{sec:domain_shift}
% ------------------------------------------------------------------------------
\subsection{Evaluation with Unknown Persons}
\label{subsec:norm_subject}
How well a model generalizes to new persons depends strongly on the evaluation
protocol. With a random split (\code{norm\_random}), images of the same person
can end up in training and test. In this thesis, groups of person and target
point are split: all frames in which a person looks at the same point are
either in training or in test, but the same person appears in both parts. The model can then use
features of this person, and generalization is overestimated. With a
subject-disjoint split (\code{norm\_subject}), the persons are split instead
(here \SI{80}{\percent} training, \SI{20}{\percent} test): all frames of a
person are only in training or only in test. Then this data leakage does not
exist (similar to leave-one-subject-out). One then measures the performance on truly
unknown persons. In the base experiments of this thesis on GazeCapture,
\code{norm\_subject} gives a test error of \SI{28.39}{\percent} and
\code{norm\_random} \SI{27.72}{\percent}. Most of this difference, however,
comes from the different test data: the constant baseline is also lower on
\code{norm\_random} (\SI{27.66}{\percent} instead of \SI{28.46}{\percent}, see
Chapter~4).

Tests across datasets are a different case, for example leave-one-dataset-out
in UniGaze~\cite{qin2026unigaze}. There, a whole dataset is held back. The two
protocols therefore measure different kinds of generalization: to unknown
persons in the same environment, or to unknown recording conditions.

\subsection{Domain Adaptation and Domain Generalization}
\label{subsec:domain_adaptation}
When training and test data differ in persons, cameras or environment,
appearance-based methods become much worse. There are two counter-strategies.
Domain adaptation uses (usually unlabelled) data from the target environment.
Domain generalization needs no target data at all.

\paragraph{CRGA.}
Wang et al.~\cite{wang2022crga} propose \emph{Contrastive Regression Gaze
Adaptation} (CRGA), an unsupervised domain adaptation method. A module for
\emph{Contrastive Domain Generalization} (CDG) learns stable features in the
source domain: features with similar gaze are pulled together, features with
very different gaze are pushed apart. A module for \emph{Contrastive
Self-training Adaptation} (CSA) then learns with pseudo-labels in the target
domain. With ETH-XGaze and Gaze360 as source and MPIIGaze, RT-GENE, GazeCapture
and EyeDiap as target, the authors improve the baseline by 40.4\,\%, 34.7\,\%,
55.8\,\% and 34.3\,\% (source ETH-XGaze) and by 31.7\,\%, 30.5\,\%, 32.9\,\%
and 23.8\,\% (source Gaze360).

CRGA therefore solves a different problem: adaptation to a new environment
without labelled target data. This thesis, in contrast, uses a few labelled
calibration data of the target person. This allows supervised training with
real labels, not only with pseudo-labels. Self-training with pseudo-labels can
also create a confirmation bias: wrong predictions are reinforced in further
training. In addition, this thesis is about differences between persons and
not about switching between benchmark datasets. CRGA is therefore related work
on robustness, but not a directly comparable baseline.

\paragraph{Domain generalization.}
Methods such as PureGaze~\cite{cheng2022puregaze} or
CLIP-Gaze~\cite{yin2024clipgaze} (\cref{subsec:clip}) try to remove
gaze-irrelevant factors from the features already during training. Then the
model also works in new environments without target data.
UniGaze~\cite{qin2026unigaze} shows that large gaze-specific pretraining can
beat these methods in tests across datasets (\cref{subsec:unigaze}).

\subsection{Gaze Estimation for Several Persons}
\label{subsec:gazeonce}
Zhang et al.~\cite{zhang2022gazeonce} present GazeOnce, a one-stage method. It
estimates the gaze of all persons in an image in a single pass and stays
real-time even with more than ten faces. Besides the gaze direction, it also
predicts face boxes and facial landmarks. A self-supervised loss checks whether
the estimated 3D direction fits its 2D projections. For training, the authors
create the synthetic dataset MPSGaze: they insert faces from ETH-XGaze with
known gaze into matching faces from WiderFace. According to the authors,
GazeOnce is more accurate and faster than existing methods, especially with
many faces.

GazeOnce therefore solves a different problem: the gaze of many persons in
free scenes, not personalization for one webcam user. Methodically interesting
is the consistency between the 3D direction and its 2D projection. MPSGaze also
shows that synthetic training data can differ from real webcam images.

\subsection{Differences in Devices and Environment}
\label{subsec:geraete}
Besides differences between datasets, many factors of device and environment
affect the accuracy: glasses and occlusions, background and light, image
resolution and camera distance, and the position of camera and screen. The
literature gives some hints:
\begin{itemize}
  \item Cheng et al.~\cite{cheng2024review} observe much larger errors at large
        distance and low resolution (Gaze360, RT-GENE).
  \item GazeCLIP~\cite{wang2026gazeclip} names bad light and occlusion by hair
        or glasses as typical failure cases.
  \item UniGaze~\cite{qin2026unigaze} stresses how important many different
        head poses and a wide range of gaze labels are.
  \item CLIP-Gaze~\cite{yin2024clipgaze} explicitly models accessories and
        image quality as gaze-irrelevant factors.
\end{itemize}
Systematic studies of these factors for personalization with short webcam
calibrations are still largely missing. This thesis therefore also looks at
two devices with different webcams, glasses, background and calibration
length across several sessions (see Chapter~4).

% ------------------------------------------------------------------------------
\section{Test-Time Adaptation}
\label{sec:tta}
% ------------------------------------------------------------------------------
In test-time adaptation (TTA), a fully trained model is adapted to the person or
the environment only during use, usually with few or no labels. Unlike classical
calibration (\cref{subsec:kalibrierung}), the user should have to do as little
as possible.

\subsection{TTP-GE: Self-Supervised and Weakly Supervised Adaptation}
\label{subsec:ttpge}
Nagpure et al.~\cite{nagpure2025ttpge} present TTP-GE (\emph{Test-Time
Personalized Gaze Estimation}), a two-stage method for gaze estimation in the
car. It is based on GazeDPTR~V2. This architecture has two branches: one for
the original face image and one for a normalized face image. A transformer
combines the features. With position information from a point projection, the
model can also recognize gaze zones.

Meta-training uses two loops based on MAML. In the inner loop, a
self-supervised consistency loss (in the style of BYOL++) adapts the model to
new persons -- without gaze labels. The outer loop optimizes gaze regression
and zone recognition with labelled data. At test time there are two steps:
first, the model adapts in a self-supervised way to unlabelled images of the
person; this helps with changed image conditions. Then the zone recognition is
adapted with a few zone labels (for example mirror, dashboard, windscreen).
These labels are collected in a short setup phase.

On IVGaze (\num{44705} images of 125 persons), the combination reaches a mean
error of $6.00^\circ$. The baseline is at $6.78^\circ$, \emph{weakly supervised
adaptation} alone (here with a few coarse zone labels) at $6.13^\circ$ and
\emph{self-supervised adaptation} alone (without any labels) at $6.72^\circ$. The share of estimates with less than $6^\circ$ error rises from
49.95\,\% to 58.2\,\%. So the two steps complement each other: self-supervision
handles the image conditions, and the zone labels provide personal
information.

\subsection{Other Label-Saving Methods}
\label{subsec:tta_weitere}
Some recent works pursue similar goals with other means:
\begin{itemize}
  \item \textbf{TPGaze}~\cite{liu2024tpgaze} adapts only a meta-learned prompt
        at test time. It has less than 1\,\% of the parameters of a ResNet-18.
        A simple symmetry loss is enough as a goal without labels. According to
        the authors, the adaptation is up to ten times faster than other
        methods. It is tested in four scenarios across datasets.
  \item \textbf{ELF-UA}~\cite{wu2024elfua} treats adaptation \emph{without
        labels} as its own problem: a few unlabelled images of the person are
        enough. Meta-learning uses a bound from domain adaptation as the loss
        function.
  \item \textbf{PIGaze}~\cite{wang2026pigaze} obtains calibration data in the
        car automatically: the gaze at known zones such as mirror and
        windscreen must fit together geometrically (self-calibration). A
        pretrained base model is a \emph{frozen} feature extractor. On top, a
        small adapter learns the personal mapping. On IVGaze and LBW, PIGaze is
        better than fine-tuning.
\end{itemize}

\paragraph{Limitations and differences.}
These methods want to reduce calibration effort with unlabelled or weakly
labelled data. They mostly estimate 3D gaze directions and, for TTP-GE and
PIGaze, work in the car with fixed zones. Their adaptation (meta-training,
online updates, EMA target networks) is much more complex than simple head
training. General pretrained models such as DINOv2, CLIP or MAE are not studied
as backbones. This thesis follows the complementary path of \emph{supervised}
personalization with screen coordinates from a short calibration. It is
interesting, however, that PIGaze -- like this thesis -- combines a frozen
backbone with a small, personally trained module.

\begin{table}[htbp]
  \centering
  \small
  \caption{Comparison of TTP-GE and this thesis.}
  \label{tab:ttpge_vergleich}
  \begin{tabularx}{\textwidth}{L{3.0cm} Y Y}
    \toprule
    \textbf{Aspect} & \textbf{TTP-GE~\cite{nagpure2025ttpge}} & \textbf{This thesis} \\
    \midrule
    Approach & Test-time adaptation with unlabelled images and a few zone labels & Supervised personalization with labelled calibration data \\
    Backbone & GazeDPTR~V2: two ResNet-18 branches with transformer fusion & ResNet-18 (ImageNet), frozen \\
    Labels & Unlabelled images and nine gaze zones & Screen coordinates; \code{WAIT} frames are removed \\
    Adaptation & Meta-training, online updates, EMA target network & Regression head only \\
    Application & Driver monitoring in the car & Webcam in front of a screen \\
    Result & $6.00^\circ$ on IVGaze & Diagonal error in \%, not directly comparable \\
    \bottomrule
  \end{tabularx}
\end{table}

% ------------------------------------------------------------------------------
\section{Data Quality, Wrong Labels and Timing}
\label{sec:label_noise}
% ------------------------------------------------------------------------------
Supervised personalization assumes that every calibration frame has the right
label, that is, the point the person really looks at in this moment. In
practice this is only roughly true. Differences between label and real gaze are
called \emph{label noise}. In short calibrations it is especially harmful,
because every wrong frame is a large part of the training data.

\subsection{Data Quality in Eye Tracking}
\label{subsec:datenqualitaet}
In eye-tracking research, data quality is usually described by three measures.
Niehorster et al.~\cite{niehorster2026quality} give a recent overview of how to
measure them and provide tools for it. For $N$ measurements in which the person
looks at a known target:
\begin{itemize}
  \item \textbf{Accuracy}: the \emph{systematic} deviation between measured
        and true gaze point, usually as mean angular error
        \[
          \mathrm{Acc} = \frac{1}{N}\sum_{i=1}^{N} \theta_i ,
        \]
        where $\theta_i$ is the angle between measured and true gaze.
  \item \textbf{Precision}: the \emph{spread} of repeated measurements,
        independent of the true target. A common measure is the RMS value
        between consecutive measurements
        \[
          \mathrm{RMS\text{-}S2S} = \sqrt{\frac{1}{N-1}\sum_{i=1}^{N-1} \theta_{i,i+1}^{2}} ,
        \]
        with $\theta_{i,i+1}$ as the angle between measurement $i$ and $i+1$,
        or the standard deviation of the measurements around their mean.
  \item \textbf{Data loss}: the share of measurements without a valid gaze
        value, $\mathrm{DL} = N_{\text{invalid}} / N \cdot \SI{100}{\percent}$.
\end{itemize}
The difference: an eye tracker can be precise but inaccurate (the measurements
lie close together but next to the target), or accurate but imprecise (the
mean is right, but the measurements spread a lot). The diagonal error of this
thesis is an RMSE and therefore contains both: the systematic deviation and the
spread. For learning methods, a
fourth measure is added: the quality of the \emph{training labels}. Wrong
labels distort not only the evaluation but also the learned model.

\subsection{Reaction Time and Saccades}
\label{subsec:latenz}
When a calibration target appears at a new place, the gaze does not follow at
once. There is a short reaction time before the \emph{saccade} (a very fast eye
movement to a new target) starts.
Valliappan et al.~\cite{valliappan2020smartphone} measure on smartphones a mean
of \SI{210}{\milli\second} (median \SI{167}{\milli\second}). This fits earlier
studies with \SIrange{200}{250}{\milli\second}. Then follow the saccade itself
and possibly small corrections until the eye rests on the target.

Frames from this \emph{transition period} have the label of the \emph{new}
position. The gaze,
however, is still on the \emph{old} position or somewhere in between. At
\SI{30}{fps}, a reaction time of about \SI{200}{\milli\second} alone affects
about six frames per jump. The more often the target jumps, the more frames
have a wrong label.

\subsection{Filtering During Recording}
\label{subsec:phase_logging}
A simple counter-measure: discard the frames directly after a target change.
Valliappan et al.~\cite{valliappan2020smartphone}, for example, always discard
the first \SI{800}{\milli\second} after each new target. A fixed time window is
easy to implement. But it is based on an assumption about the length of the
transition, and its effect is usually not studied on its own.

This thesis takes a similar but more precise path. The recording program stores
the current \emph{phase} of the target for every frame: \code{WAIT} (a short
waiting time of \SI{0.5}{\second} directly after a jump) or \code{MOVE} (the
target then moves for \SI{4.5}{\second}). This makes it possible to identify
the \code{WAIT} frames exactly and reproducibly and to remove them before
training. In the recordings, this affects about 10\,\% of the frames. Unlike a
fixed time window, this logging allows a direct comparison: models with and
without cleaned data under otherwise equal conditions.

\subsection{Filtering After Recording and Robust Learning}
\label{subsec:dmrf}
Other methods detect wrong labels only \emph{after} recording, with statistical
or model-based criteria.

\paragraph{DMRF.}
Jiang et al.~\cite{jiang2026dmrf} propose \emph{Dynamic Multi-model Regression
Filtering} (DMRF), a general method against wrong labels in regression tasks.
It solves a \emph{chicken-and-egg problem} (two things depend on each other,
so it is unclear where to start): wrong labels distort the models used to
estimate the noise, and distorted models estimate the noise wrongly. DMRF has
three parts:
\begin{enumerate}
  \item a \emph{warm-up filtering} with a distribution function of the
        estimated errors,
  \item a \emph{dynamic model selection} based on a correlation criterion,
  \item a \emph{data recall} step that brings examples with small loss back to
        the clean data.
\end{enumerate}
Besides 15 benchmarks from UCI and KEEL and an age estimation task, DMRF is
built into the gaze method L2CS-Net. On a pre-processed MPIIGaze dataset with
\num{45000} images (leave-one-person-out), the mean error drops from
$3.92^\circ$ to $3.85^\circ$, while 74.3\,\% of the data are kept. The code is
public.

\paragraph{SeeTN.}
Peng et al.~\cite{peng2026seetn} are the first to study systematically how
wrong labels affect domain generalization in gaze estimation. In a pilot
experiment they corrupt 10, 20 and 30\,\% of the Gaze360 labels with strong
Gaussian noise. The error of a ResNet-18 baseline on MPIIGaze then rises from
$7.94^\circ$ to $12.17^\circ$ (20\,\%) and $14.38^\circ$ (30\,\%). Their method
\emph{See-Through-Noise} builds a semantic feature space. It detects wrong
examples when feature and label do not fit together, and it transfers gaze
information from clean to noisy examples.

\paragraph{Filter before or after?}
Methods such as DMRF and SeeTN detect wrong labels from the data themselves.
They therefore also find unexpected sources of error. But they need enough
examples for a reliable estimate, compute more, and have so far only been tested
on large datasets, not on short calibrations of one person. Phase logging, in
contrast, only covers one \emph{known} source of error that comes from the
recording procedure -- but it does so exactly and without an extra model. The
two approaches complement each other (\cref{tab:dmrf_vergleich}).

\begin{table}[htbp]
  \centering
  \small
  \caption{Comparison of filtering after recording (DMRF) and phase logging.}
  \label{tab:dmrf_vergleich}
  \begin{tabularx}{\textwidth}{L{3.0cm} Y Y}
    \toprule
    \textbf{Criterion} & \textbf{DMRF~\cite{jiang2026dmrf}} & \textbf{This thesis} \\
    \midrule
    Strategy & After: estimate errors and filter & Before: store the phase during recording \\
    Detection & Model predictions, error patterns, dynamic model selection & \code{WAIT}/\code{MOVE} label per frame \\
    Errors covered & All statistically unusual examples & Only the transition after a jump \\
    Data kept & 74.3\,\% (MPIIGaze) & Approx.\ 90\,\% \\
    Application & Many persons, leave-one-person-out & Webcam calibration of one person \\
    Extra effort & Several regression models, iterative training & No extra model \\
    Result & $3.92^\circ \rightarrow 3.85^\circ$ & Diagonal error in \%, not directly comparable \\
    \bottomrule
  \end{tabularx}
\end{table}

% ------------------------------------------------------------------------------
\section{Comparison and Critical Discussion}
\label{sec:vergleich}
% ------------------------------------------------------------------------------
\Cref{tab:sota_vergleich} summarizes the most important works of this chapter.
The columns are: goal, features used, type of personalization and handling of
label quality.

\begin{table}[htbp]
  \centering
  \footnotesize
  \setlength{\tabcolsep}{3pt}
  \caption[Comparison of the most important works]{Comparison of the most important works. $k$ = number of calibration images per person. The errors come from different protocols and metrics and are not directly comparable.}
  \label{tab:sota_vergleich}
  \begin{tabularx}{\textwidth}{L{2.3cm} L{2.4cm} L{2.6cm} L{2.6cm} Y L{1.9cm}}
    \toprule
    \textbf{Work} & \textbf{Goal} & \textbf{Features} & \textbf{Personalization} & \textbf{Label quality} & \textbf{Result} \\
    \midrule
    iTracker~\cite{krafka2016itracker} & 2D gaze, mobile & CNN branches similar to AlexNet & SVR, up to 13 points & -- & \SI{1.34}{\centi\metre} (smartphone) \\
    FAZE~\cite{park2019faze} & Few images & DenseNet DT-ED, supervised & MAML, $k \le 256$ & -- & $3.14^\circ$ ($k=9$) \\
    SPAZE~\cite{linden2019spaze} & Few images & CNN, latent person parameters & Small optimization problem & -- & $2.70^\circ$ ($k=9$) \\
    Chen \& Shi~\cite{chen2020offset} & Offset calibration & CNN & Offset from one gaze point & -- & up to 10\,\% better \\
    Valliappan~\cite{valliappan2020smartphone} & Smartphone & CNN, pretrained on GazeCapture & Fine-tuning + regression, approx.\ 30\,s & Fixed 800-ms discard & \SI{0.46}{\centi\metre} \\
    CRGA~\cite{wang2022crga} & Unsupervised DA & ResNet (ImageNet) & Pseudo-labels & -- & up to 55.8\,\% better \\
    UniGaze~\cite{qin2026unigaze} & Generalization & MAE-ViT, 1.6~million faces & Full fine-tuning & -- & $4.65$--$5.86^\circ$ (DG) \\
    CLIP-Gaze~\cite{yin2024clipgaze} & Generalization & CLIP (frozen) + ResNet-18 & -- & Gaze-irrelevant factors & $6.97^\circ$ (DG) \\
    GazeCLIP~\cite{wang2026gazeclip} & Supervised regression & CLIP ResNet-50 (fine-tuned) & -- & -- & $3.5^\circ$ (MPIIFaceGaze) \\
    TTP-GE~\cite{nagpure2025ttpge} & Test-time adaptation & GazeDPTR~V2 & Self- + weakly supervised & -- & $6.00^\circ$ (IVGaze) \\
    TPGaze~\cite{liu2024tpgaze} & Test-time adaptation & ResNet-18 & Prompt (<1\,\% of parameters) & -- & -- \\
    DMRF~\cite{jiang2026dmrf} & Filter wrong labels & L2CS-Net & -- & After, model-based & $3.85^\circ$ (MPIIGaze) \\
    SeeTN~\cite{peng2026seetn} & Robust generalization & ResNet-18/50 & -- & After, feature-based & -- \\
    \midrule
    \textbf{This thesis} & Webcam personalization & ResNet-18 (ImageNet), frozen & Head only; \SIrange{30}{120}{\second} & Before, phase logging & Diagonal error in \% \\
    \bottomrule
  \end{tabularx}
\end{table}

\paragraph{Critical discussion.}
The comparison leads to four observations:
\begin{enumerate}
  \item \textbf{Features and personalization are studied separately.} Works on
        large pretrained models (UniGaze, CLIP-Gaze, GazeCLIP) aim at
        generalization across datasets and use full fine-tuning with a lot of
        data. Works on personalization (FAZE, SPAZE, Chen \& Shi), in contrast,
        use CNNs that were trained specially for gaze. How well general
        pretrained models can be personalized with little calibration data is
        still open.
  \item \textbf{Few parameters are often enough.} Offset calibration, few person
        parameters, prompt tuning and frozen models with adapters (PIGaze) all
        show: a large part of the personal error can be explained with a few
        adapted parameters. Kumar et al.~\cite{kumar2022finetuning} give a
        general reason for this.
  \item \textbf{Wrong labels are usually handled afterwards.} DMRF and SeeTN
        filter or weight wrong labels in large datasets. Errors caused by the
        procedure of short calibrations are handled, if at all, with fixed time
        windows (Valliappan et al.), without studying their effect in detail.
  \item \textbf{The results are hard to compare.} The works report angles in
        degrees, distances in centimetres or relative improvements, and they use
        different datasets, splits and protocols. An angular error can only be
        converted into a screen error if the geometry (distance, screen size)
        is known. The metric of this thesis, the Euclidean error relative to the
        screen diagonal, is therefore not directly comparable with the
        literature. The numbers in \cref{tab:sota_vergleich} are for
        orientation, not a ranking.
\end{enumerate}

% ------------------------------------------------------------------------------
\section{Research Gaps and Link to This Thesis}
\label{sec:luecken}
% ------------------------------------------------------------------------------
\subsection{Research Gaps}
\label{subsec:luecken}
The state of the art shows four research gaps:

\begin{enumerate}[label=\textbf{G\arabic*:}, leftmargin=*, align=left]
  \item \textbf{Data-efficient personalization of pretrained models.}
  Large pretrained models (ImageNet CNNs, DINOv2, CLIP, MAE or UniGaze) have so
  far been studied for gaze estimation mainly with full fine-tuning on large
  datasets. Studies on how well they can be adapted to one person with only a
  few minutes of webcam calibration are missing.

  \item \textbf{Lightweight adaptation for personalization.}
  Several works show that a few adapted parameters capture most of the personal
  error. Still, head-only training, adapters and LoRA have hardly been compared
  systematically for personalization in gaze estimation.

  \item \textbf{Label quality of short calibrations.}
  Reaction time and target changes produce calibration frames with a wrong
  label. Existing work mostly handles wrong labels afterwards in large datasets
  or discards fixed time windows. How an exact, protocol-based removal of such
  frames affects personalization has not been studied yet.

  \item \textbf{Realistic evaluation under webcam conditions.}
  Many results come from benchmarks with controlled conditions or from random
  splits on image level, which can overestimate generalization. How camera,
  glasses, background and calibration length affect personalization is
  hardly documented.
\end{enumerate}

\subsection{What This Thesis Contributes}
\label{subsec:ueberleitung}
This thesis addresses the gaps with a simple and reproducible personalization
pipeline for webcam eye tracking:

\begin{itemize}
  \item \textbf{Frozen backbone with a small regression head (G1, G2):} A
        ResNet-18 pretrained on ImageNet stays frozen. Its features are stored,
        and only a small regression head is trained per person. I also test whether
        training more layers helps. Other backbones (DINOv2, CLIP, MAE,
        UniGaze) and adapters or LoRA are the next step (Chapter~6).
  \item \textbf{Phase logging and cleaning (G3):} The recording program stores
        the phases \code{WAIT} and \code{MOVE}. This allows a controlled
        measurement of the effect of removing transition frames: models with
        original and with cleaned data are compared under otherwise equal
        conditions.
  \item \textbf{Systematic ablations:} Batch size, learning rate, trainable
        layers, activation, loss function (including an L1 loss weighted by
        position) and amount of data are studied with repeated runs and
        confidence intervals.
  \item \textbf{Realistic evaluation (G4):} The base models are evaluated with
        the subject-disjoint split \code{norm\_subject}. Personalization is
        tested on 18 sessions with changing webcam (two devices), glasses,
        background and calibration length, with both a random and a temporal
        split.
\end{itemize}

Chapter~3 describes the methodology and Chapter~4 the results. Chapter~5 uses
them to answer the research questions from Chapter~1.

% ==============================================================================
